\documentclass[10pt,a4paper]{amsart}
\usepackage[margin=19mm]{geometry}
\usepackage{amsmath,amssymb,mathtools}
\usepackage[scr=rsfs,cal=euler]{mathalfa}
\usepackage{xfrac}
\usepackage[unicode,hidelinks,bookmarks=false]{hyperref}
\newcommand{\ie}{i.e.\ }
\newcommand{\cf}{cf.\ }
\newcommand{\resp}{respectively\ }
\newcommand{\Subsection}{\subsection}
\providecommand{\nobreakdash}{\nobreak\hbox{-}}
\setlength{\emergencystretch}{2em}
\multlinegap0pt
\usepackage{amssymb}
\usepackage{algorithm}\usepackage{algpseudocode}
\newtheorem{theorem}{Theorem}
\title{Kernel Methods for Some Transport Equations with Application to Learning Kernels for the Approximation of Koopman Eigenfunctions: \\ A Unified Approach via Variational Methods, Green’s Functions and the Method of Characteristics}
\author{Boumediene Hamzi, Houman Owhadi, Umesh Vaidya}
\date{Source: arXiv:2603.06872v1 (CC BY 4.0)}
\begin{document}
\maketitle

\noindent\emph{Source and licence.} Kernel Methods for Some Transport Equations with Application to Learning Kernels for the Approximation of Koopman Eigenfunctions: \\ A Unified Approach via Variational Methods, Green’s Functions and the Method of Characteristics. Version arXiv:2603.06872v1 (CC BY 4.0), distributed by arXiv under the Creative Commons Attribution 4.0 International licence, http://creativecommons.org/licenses/by/4.0/, as recorded in the arXiv licence metadata for this exact version and shown on the abstract page https://arxiv.org/abs/2603.06872v1. Full licence notice: \texttt{licenses/arxiv-2603.06872-CC-BY-4.0.txt}.\par\smallskip

\noindent\textit{Editorial scope.} Counted unit: the article's theorem thm:truncated-characteristic-kernel (source0.tex lines 1443--1460). The article's complete original proof of it is delivered with it (source0.tex lines 1461--1548, one \texttt{proof} environment of 88 lines). No mathematics is written, repaired or re-derived: the statement and the proof are the article's own lines, in the article's own notation, and no retained line is substituted or re-flowed.  No further source-local context is needed: every cross-reference of every retained line resolves inside the two delivered spans.  Nothing was omitted from any retained span, so the package carries no omission record.  No retained line contains a citation, so the wrapper carries no bibliography and no bibliography entry.  No retained line contains a figure environment, an image-inclusion command or drawing code, so no image is delivered and the package declares no asset dependency.  Editorial formatting only, in addition: an AMS \texttt{amsart} preamble that restates the article's own declarations and macros used by the retained lines, namely the environments \texttt{theorem} on separate counters, together with the standard packages those lines need.  The retained environments print in this document as Theorem 1 in order of appearance, whereas the article prints the same items with the numbering of its own full document.  The complete native source of the article as distributed in the arXiv e-print for this exact version is bundled unchanged as \texttt{sources/195902/source0.tex}.
\begin{theorem}[Truncated Flow Kernel with Approximate Koopman Behavior]
\label{thm:truncated-characteristic-kernel}
Let \( f \in C^1(\Omega; \mathbb{R}^d) \) generate a complete flow \( \varphi(s; x) \) on a compact domain \( \Omega \), and fix \( T > 0 \), \( \lambda \in \mathbb{C} \) with \( \Re(\lambda) > 0 \), and \( w \in \mathbb{R}^d \). Define the truncated characteristic coordinate:
\[
\xi_T(x) := w^\top x + \int_{-T}^0 e^{-\lambda(-s)} w^\top {\mathbb f}(\varphi(s; x)) \, ds.
\]
Then the kernel
\[
K_T(x, y) := \xi_T(x) \cdot \xi_T(y)
\]
is continuous, symmetric, and positive semi-definite. It defines a one-dimensional reproducing kernel Hilbert space \( \mathcal{H}_K = \operatorname{span}\{\xi_T\} \subset C^0(\Omega) \). Furthermore, \( \xi_T \) approximately satisfies the Koopman eigenvalue equation:
\[
{\mathbb f}(x) \cdot \nabla \xi_T(x) = \lambda \xi_T(x) + R_T(x),
\quad
\text{where } R_T(x) = -e^{-\lambda T} w^\top {\mathbb f}(\varphi(-T; x)).
\]
As \( T \to \infty \), the residual term \( R_T(x) \to 0 \) uniformly, and \( \xi_T(x) \to \xi(x) \), recovering the exact solution to \( f \cdot \nabla \xi = \lambda \xi \).
\end{theorem}

\begin{proof}
\textbf{Step 1: Regularity of \( \xi_T(x) \).}  
Since \( f \in C^1(\Omega) \) and \( \Omega \) is compact, the flow \( \varphi(s; x) \) is smooth in both \( s \) and \( x \), and bounded over \( s \in [-T, 0] \). The integrand
\[
e^{\lambda s} w^\top {\mathbb f}(\varphi(s; x))
\]
is smooth and bounded for each \( x \in \Omega \), and the integral over a finite interval \( [-T, 0] \) is uniformly convergent. Therefore, \( \xi_T(x) \in C^0(\Omega) \), and \( K_T(x, y) = \xi_T(x) \cdot \xi_T(y) \in C^0(\Omega \times \Omega) \).

\medskip

\textbf{Step 2: Symmetry and positive semi-definiteness.}  
The kernel is clearly symmetric since \( K_T(x, y) = \xi_T(x) \cdot \xi_T(y) = K_T(y, x) \).  

To verify positive semi-definiteness, take any finite set \( \{x_i\}_{i=1}^n \subset \Omega \) and coefficients \( \{c_i\} \subset \mathbb{R} \). Then:
\[
\sum_{i,j=1}^n c_i c_j K_T(x_i, x_j)
= \sum_{i,j=1}^n c_i c_j \xi_T(x_i) \xi_T(x_j)
= \left( \sum_{i=1}^n c_i \xi_T(x_i) \right)^2 \ge 0.
\]

\medskip

\textbf{Step 3: RKHS structure.}  
As with the infinite-time kernel, the kernel \( K_T \) is rank-one. The corresponding RKHS is:
\[
\mathcal{H}_{K_T} = \left\{ \phi : \phi(x) = c \xi_T(x), \; c \in \mathbb{R} \right\},
\]
with inner product defined by:
\[
\langle \phi_1, \phi_2 \rangle_{K_T} = \frac{c_1 c_2}{\| \xi_T \|^2_{L^2(\Omega)}}.
\]

\medskip

\textbf{Step 4: Approximate Koopman property.}  
We compute \( {\mathbb f}(x) \cdot \nabla \xi_T(x) \). Define:
\[
\tilde{\xi}_T(x) := \int_{-T}^0 e^{\lambda s} w^\top {\mathbb f}(\varphi(s; x)) \, ds,
\quad \text{so that } \xi_T(x) = w^\top x + \tilde{\xi}_T(x).
\]
To compute the directional derivative along \( f \), we consider:
\[
{\mathbb f}(x) \cdot \nabla \tilde{\xi}_T(x) = \frac{d}{dt} \bigg|_{t=0} \tilde{\xi}_T(\varphi(t; x)).
\]

By the flow composition property \( \varphi(s; \varphi(t; x)) = \varphi(s + t; x) \), we get:
\[
\tilde{\xi}_T(\varphi(t; x)) = \int_{-T}^0 e^{\lambda s} w^\top {\mathbb f}(\varphi(s + t; x)) \, ds
= \int_{-T + t}^{t} e^{\lambda (r - t)} w^\top {\mathbb f}(\varphi(r; x)) \, dr,
\]
where we changed variables \( r = s + t \).

Thus,
\[
\tilde{\xi}_T(\varphi(t; x)) = e^{-\lambda t} \int_{-T + t}^{t} e^{\lambda r} w^\top {\mathbb f}(\varphi(r; x)) \, dr.
\]

Differentiating with respect to \( t \) at \( t = 0 \), we obtain:
\[
\frac{d}{dt} \tilde{\xi}_T(\varphi(t; x)) \bigg|_{t = 0}
= -\lambda \tilde{\xi}_T(x) + w^\top {\mathbb f}(x) - e^{-\lambda T} w^\top {\mathbb f}(\varphi(-T; x)).
\]

Therefore:
\[
{\mathbb f}(x) \cdot \nabla \xi_T(x) = w^\top {\mathbb f}(x) + {\mathbb f}(x) \cdot \nabla \tilde{\xi}_T(x)
= \lambda \tilde{\xi}_T(x) + w^\top {\mathbb f}(x) - e^{-\lambda T} w^\top {\mathbb f}(\varphi(-T; x)).
\]

Using \( \xi_T(x) = w^\top x + \tilde{\xi}_T(x) \), we obtain:
\[
{\mathbb f}(x) \cdot \nabla \xi_T(x) = \lambda \xi_T(x) + R_T(x),
\]
where the residual is:
\[
R_T(x) := -e^{-\lambda T} w^\top {\mathbb f}(\varphi(-T; x)).
\]

\medskip

\textbf{Step 5: Convergence of the residual.}  
Since \( f \) is continuous and \( \Omega \) compact, \( w^\top {\mathbb f}(\varphi(-T; x)) \) is uniformly bounded for \( x \in \Omega \), say by \( M > 0 \). Thus:
\[
|R_T(x)| \le M e^{-\Re(\lambda) T} \to 0 \quad \text{uniformly as } T \to \infty.
\]

Hence, as \( T \to \infty \), \( \xi_T(x) \to \xi(x) \) uniformly, and the Koopman residual vanishes uniformly. This completes the proof.
\end{proof}

\end{document}
